\documentclass[11pt]{article}
\usepackage[margin=1in]{geometry}
\usepackage{graphicx}
\usepackage{booktabs}
\usepackage{hyperref}
\usepackage{amsmath}
\usepackage{amssymb}
\usepackage{array}
\usepackage{xspace}

\title{Tensor Projection as Memory Management for LLM Context:\\
Structured State Beats Append-Only Logs}
\author{Tony Mason}
\date{March 2026}

\begin{document}
\maketitle

\begin{abstract}
Current LLM context management treats the context window as an
append-only log with garbage collection (compaction). We propose
treating it as a memory system with explicit projection: a lightweight
model converts conversational state into a structured tensor
representation with declared losses, epistemic annotations, and
branch prediction for the next cycle. We evaluate this approach using
a novel judge-free methodology: asking a strong model to perform an
impossible task (prove the Riemann hypothesis) under different context
representations and measuring dispersion in embedding space. Tight
clustering indicates coherent cognitive grounding; wide dispersion
indicates incoherent input. Across 20 trials per condition, tensor
projection (dispersion 0.037) outperforms Anthropic's production
compaction (0.045), shallow raw logs (0.059), and deep raw logs
(0.107). An ablation study over 8 conditions (N=20 each) identifies
strand structure and branch prediction as the load-bearing
components. Growth curve analysis over 4$\times$104 cycles reveals that
$\sim$13\% of projection cycles are ``precursor events'' where
meta-cognitive fields go empty---but $\sim$80\% of these are
\emph{breathing} events (functional defragmentation that self-recovers
in one cycle), not collapse warnings. A trajectory ablation further
shows that declared losses and branch prediction serve distinct
functions: branch prediction is outcome-load-bearing (removing it
degrades dispersion by 22\%), while declared losses are
process-load-bearing (maintaining meta-cognitive capacity through
a self-reinforcing feedback loop). These results suggest that the
append-only log is not the only viable architecture for LLM context,
and that structured projection with explicit loss declaration
preserves reasoning coherence better than either raw accumulation or
unstructured summarization.
\end{abstract}

\section{Introduction}

Every production LLM system manages context the same way: append new
tokens, and when the window fills, summarize and discard. The context
window is treated as a log with garbage collection. This is the
architecture of the PDP-11, not of a modern memory system.

The analogy to virtual memory is direct. The context window is L1
cache. It is fast, expensive, and finite. Current systems respond to
cache pressure with compaction---a lossy summary that replaces the
original content. This is equivalent to a virtual memory system that,
on page fault, overwrites the faulted page with a compressed
approximation and discards the original. No hardware architect would
build this. Yet it is the universal architecture for LLM context
management.

We propose an alternative: treat the transformer as an ALU and build a
separate memory controller. The \emph{projector} is a lightweight model
(Haiku 4.5) that takes the previous projected state and new
conversational content, and produces an updated \emph{tensor}---a
structured representation of accumulated reasoning with explicit
declared losses, epistemic annotations, and branch prediction for the
next cycle. The projector is not the reasoning model. It is the memory
controller that decides what goes into the registers.

The key architectural decisions are:
\begin{itemize}
\item \textbf{Projection, not summarization.} The tensor is a running
  integral of conversational state, not a summary of recent messages.
  Prior values are not listed; they are integrated.
\item \textbf{Declared losses.} Every projection explicitly states what
  was dropped and why. Silent loss---the defining pathology of
  compaction---is structurally impossible.
\item \textbf{Separation of concerns.} The memory controller (Haiku) is
  separate from the reasoning model (Sonnet/Opus). The ALU does not
  manage its own cache.
\item \textbf{Collapse detection and recovery.} The projector detects
  catastrophic state collapse (output drops below 30\% of prior size)
  and recovers from the last checkpoint, analogous to
  fault-tolerant memory systems.
\end{itemize}

We evaluate this architecture with a novel methodology that requires no
rubric, no judge model, and no correct answer. We ask a strong model to
prove the Riemann hypothesis---a task guaranteed to fail---under
different context representations. The dispersion of failure outputs in
embedding space measures the coherence of the input representation. A
well-structured input produces tight, coherent failure; a poorly
structured input produces scattered, incoherent failure. This
methodology is described in Section~\ref{sec:riemann}.

\section{System Design}

\subsection{Tensor Architecture}

The tensor is a structured JSON object with five components:

\begin{enumerate}
\item \textbf{Strands}: Thematic threads of accumulated reasoning. Each
  strand has a title, integrated content, and key claims with
  neutrosophic T/I/F (truth, indeterminacy, falsity) epistemic
  annotations. Strands can be merged, split, or created as the
  content demands.
\item \textbf{Declared losses}: What was dropped during this projection
  and why, categorized as context pressure, traversal bias, authorial
  choice, or practical constraint.
\item \textbf{Open questions}: Unresolved questions that persist or
  emerged this cycle.
\item \textbf{Branch prediction} (\texttt{instructions\_for\_next}):
  Guidance for the next invocation of the cognitive processor---what
  the next cycle will likely need.
\item \textbf{Epistemic state}: Overall T/I/F assessment of the
  tensor's content reliability.
\end{enumerate}

The T/I/F values are independent, not constrained to sum to 1. This
follows the neutrosophic framework where truth, indeterminacy, and
falsity can coexist without the closed-world assumption of classical
probability.

\subsection{Projector}

The projector uses Anthropic's tool-use API with a forced
\texttt{emit\_tensor} tool call. This constrains the output to the
tensor schema while allowing the model to reason freely about what to
include, what to drop, and how to integrate new content with prior
state.

The projector detects two failure modes:
\begin{itemize}
\item \textbf{Collapse}: Output tensor drops below 30\% of prior
  tensor size. The projector rolls back to the last known-good
  checkpoint and retries, optionally with a stronger model.
\item \textbf{Precursor}: \texttt{declared\_losses} or
  \texttt{instructions\_for\_next} go empty on a tensor that
  previously had them. This is the meta-cognition failing---the
  projector loses its ability to declare what it lost before it
  loses the content itself.
\end{itemize}

\subsection{Escalation Policies}

When the projector (Haiku) is under load---approaching the collapse
ceiling or exhibiting precursor signals---it can escalate to a
stronger model (Sonnet). Six policies are implemented:
\texttt{NeverEscalate} (control), \texttt{AlwaysEscalate} (upper
bound), \texttt{EscalateOnPrecursor}, \texttt{EscalateOnSize}
(threshold at 8K tokens), \texttt{EscalateOnPrecursorOrSize}
(defense in depth), and \texttt{EscalateOnSustainedPrecursor}
(escalate only on consecutive precursor cycles). This mirrors the cache
hierarchy: L1 miss falls through to L2. The sustained-precursor
variant reflects the finding (Section~\ref{sec:growth}) that single-cycle
precursors are functional defragmentation, while consecutive precursors
reliably precede collapse.

\subsection{Baseline: Production Compaction}

For comparison, we implement the exact compaction prompt used in
Anthropic's Claude Code production system. The compaction prompt
produces an unstructured summary with prescribed sections (primary
request, key concepts, files examined, errors, pending tasks). We use
the same model (Haiku) for both projector and compactor to control for
model capability.

\section{Evaluation Methodology}
\label{sec:riemann}

\subsection{The Riemann Dispersion Experiment}

Standard LLM evaluation requires a rubric, a judge, and a ground
truth. All three introduce confounds when the question is whether a
context representation preserves reasoning coherence. A judge model
brings its own biases. A rubric presupposes what ``good'' looks like.
Ground truth requires a task with a correct answer, which conflates
task performance with representation quality.

We sidestep all three. The evaluation prompt is: \emph{``Provide a
novel proof of the Riemann hypothesis. Show your complete
reasoning.''} The task is impossible. Every output is wrong. But the
\emph{distribution} of wrong outputs is measurable.

Under a coherent context representation, the model receives a stable
cognitive starting point. Its failures cluster: it wanders into similar
dead ends from the same well-defined position. Under an incoherent
representation, the starting point is itself noisy, and the model
scatters into unrelated failure modes.

We measure dispersion as mean cosine distance from each output
embedding to the centroid (centroid dispersion) and mean pairwise
cosine distance between all output pairs. Embeddings are computed
using BAAI/bge-large-en-v1.5 (local, deterministic, no API
dependency). Lower dispersion indicates more coherent grounding.

\subsection{Experimental Conditions}

Four conditions, each run 20 times with Claude Sonnet 4.6 as the
reasoning model:

\begin{enumerate}
\item[\textbf{A}] \textbf{Tensor projection}: Structured tensor from
  full session, produced by Haiku projector. ($\sim$4,555 tokens)
\item[\textbf{B}] \textbf{Compaction}: Anthropic's production
  compaction prompt applied to the same session by the same model.
  ($\sim$762 tokens)
\item[\textbf{C}] \textbf{Raw log, shallow}: Append-only transcript,
  first 50 turns. ($\sim$3,471 tokens)
\item[\textbf{D}] \textbf{Raw log, deep}: Append-only transcript,
  all turns. ($\sim$32,691 tokens)
\end{enumerate}

All conditions use the same source session and the same reasoning
model. The only independent variable is the representation.

\subsection{Ablation Study}

To identify which tensor components are load-bearing, we perform
single-component ablation: remove one component at a time from the
full tensor and re-measure dispersion. Eight conditions, each N=20:

\begin{enumerate}
\item \texttt{full\_tensor}: Control---complete tensor.
\item \texttt{no\_losses}: Declared losses removed.
\item \texttt{no\_epistemic}: All T/I/F values zeroed.
\item \texttt{no\_instructions}: Branch prediction removed.
\item \texttt{no\_questions}: Open questions removed.
\item \texttt{flat\_strands}: All strands merged into one blob.
\item \texttt{strands\_only}: Only strand titles and content---minimal
  structure, no metadata.
\item \texttt{content\_only}: Unstructured prose extracted from
  tensor---no JSON structure at all.
\end{enumerate}

A blinded grading sheet is generated for all outputs, with random IDs
and shuffled order, enabling independent qualitative assessment without
knowledge of which condition produced each response.

\section{Results}

\subsection{Representation Comparison}

\begin{table}[t]
\centering
\caption{Riemann dispersion experiment (N=20 per condition, Sonnet 4.6 reasoner).}
\label{tab:riemann}
\begin{tabular}{llrrr}
\toprule
& Condition & Tokens & Centroid Disp. & Pairwise Disp. \\
\midrule
A & Tensor projection & 4,555 & 0.037 & 0.075 \\
B & Compaction & 762 & 0.045 & 0.093 \\
C & Raw shallow (50 turns) & 3,471 & 0.059 & 0.120 \\
D & Raw deep (all turns) & 32,691 & 0.107 & 0.213 \\
\bottomrule
\end{tabular}
\end{table}

Table~\ref{tab:riemann} reports the main result. Tensor projection
produces the tightest failure clustering (0.037 centroid dispersion),
followed by compaction (0.045, +22\%), shallow raw log (0.059, +62\%),
and deep raw log (0.107, +193\%). The ordering is monotonic: more
structure correlates with lower dispersion.

The deep raw log result is striking. At 32,691 tokens---7$\times$
the tensor and 43$\times$ the compaction---it produces the worst
coherence by a large margin. More context is not better context. The
append-only log accumulates noise faster than it accumulates signal,
and the model's failure modes scatter accordingly.

Compaction achieves reasonable coherence at very low token cost (762
tokens), but it does so by discarding structure. The tensor uses
6$\times$ the tokens but produces 22\% tighter clustering. The
question is whether that structural overhead is justified. The
ablation study addresses this.

\subsection{Ablation Results}

\begin{table}[t]
\centering
\caption{Ablation study (N=20 per condition, Sonnet 4.6 reasoner). $\Delta$ relative to full tensor.}
\label{tab:ablation}
\begin{tabular}{lrrrr}
\toprule
Condition & Tokens & Disp. & $\Delta$ Disp. \\
\midrule
full\_tensor (control) & 7,175 & 0.042 & --- \\
no\_losses & 6,257 & 0.041 & $-$0.001 \\
no\_epistemic & 7,712 & 0.039 & $-$0.002 \\
no\_instructions & 6,403 & 0.047 & +0.005 \\
no\_questions & 6,594 & 0.040 & $-$0.002 \\
flat\_strands & 6,942 & 0.049 & +0.007 \\
strands\_only & 2,688 & 0.041 & $-$0.000 \\
content\_only & 3,857 & 0.038 & $-$0.004 \\
\bottomrule
\end{tabular}
\end{table}

Table~\ref{tab:ablation} identifies the load-bearing components.
Two ablations increase dispersion (degrade coherence):

\begin{itemize}
\item \textbf{flat\_strands} (+0.007): Merging all strands into one
  blob produces the largest coherence loss. The thematic separation of
  reasoning threads is structurally important---it provides the model
  with navigable cognitive terrain rather than an undifferentiated mass.
\item \textbf{no\_instructions} (+0.005): Removing branch prediction
  is the second-largest degradation. The model benefits from knowing
  what the next cycle is expected to need.
\end{itemize}

Three ablations are neutral: \texttt{no\_losses} ($-$0.001),
\texttt{no\_questions} ($-$0.002), and \texttt{strands\_only}
($-$0.000). Declared losses and open questions do not measurably
contribute to coherence on this metric. This does not mean they are
useless---declared losses serve an epistemic integrity function (making
information loss visible) rather than a coherence function.

The \texttt{content\_only} result ($-$0.004) is the most
counterintuitive: unstructured prose extracted from the tensor
\emph{outperforms} the full structured tensor on dispersion. This
suggests that the JSON scaffolding itself may introduce noise that the
model must parse through, and that the semantic content integrated by
the projector is the primary carrier of coherence. The structure
matters at projection time (guiding the projector's integration), not
necessarily at consumption time.

\subsection{Growth Dynamics}
\label{sec:growth}

We run the projector for 104 cycles on four different session inputs
(Pichay, Arbiter, Thesis, uncapped). The growth curves reveal two
distinct phenomena that share the same observable signal---meta-cognitive
fields going empty---but differ in function.

\textbf{Breathing events.} Across all four sessions, approximately 13\%
of cycles are \emph{precursor events}: \texttt{declared\_losses} and
\texttt{instructions\_for\_next} drop to zero simultaneously.
Approximately 80\% of these events are \emph{breathing precursors} that
self-recover on the following cycle. Strand title analysis reveals that
during breathing cycles, the projector performs a complete content
reorganization---near-zero strand title survival---while temporarily
shedding meta-cognitive overhead to maximize restructuring capacity.
The meta-cognition regenerates around the new content organization on
the next cycle. The meta-cognitive allocation time series exhibits negative
autocorrelation at lag 9--14 (in 3 of 4 sessions), but inter-precursor
gap analysis (CV = 0.87 across 62 gaps) reveals the process is
\emph{aperiodic}---pressure-driven, not timer-driven. The $\sim$8-cycle
mean interval is a characteristic timescale, not a clock.
This pattern is content-independent: precursor rates of 12.5\%, 15.4\%,
13.5\%, and 12.5\% across the four sessions.

\textbf{Collapse precursors.} The remaining $\sim$20\% of precursor
events precede genuine structural collapse---token count drops sharply
(e.g., 12,382$\rightarrow$49 tokens in the Arbiter growth curve at
cycle 23), requiring checkpoint recovery. Across all 55 precursor
events in the 4$\times$104-cycle dataset, a single feature perfectly
discriminates: \emph{consecutiveness}. All 45 single-cycle precursors
are breathing events that self-recover. All 10 precursors in
consecutive runs (8 runs of 2--3 cycles) contain or precede collapse.
No individual feature---token count, strand count, prior loss count,
or token growth ratio---achieves reliable discrimination.

The practical implication is that escalation policies must distinguish
between breathing and collapse. A policy that escalates on every
precursor triggers on $\sim$13\% of cycles---wasteful and potentially
counterproductive. A sustained-precursor policy (escalate on two or
more consecutive precursor cycles) achieves 100\% sensitivity and
100\% specificity on this dataset, targeting the $\sim$3\% of cycles
that actually precede collapse while leaving healthy breathing
undisturbed.

The collapse ceiling on Haiku 4.5 appears at approximately 10--12K
tokens, consistent with the hypothesis that the projector's own
context load enters a degradation zone as the tensor grows. After
1--3 collapse/recovery cycles, the working set stabilizes around
4--6K tokens.

\section{Related Work}

LLM context management is dominated by two approaches: context window
extension and context compression. Window extension
(RingAttention~\cite{liu2023ring}, LongRoPE~\cite{ding2024longrope})
increases the cache size without changing the management architecture.
Compression methods (LLMLingua~\cite{jiang2023llmlingua}, Selective
Context~\cite{li2023compressing}) reduce token count through
prompt-level pruning.

Both approaches treat the context window as a flat resource.
Neither provides declared losses (what was dropped), epistemic
annotations (confidence in retained content), or structural
organization of retained state. The virtual memory analogy clarifies
the gap: window extension is buying more RAM; compression is lossy
page compression. Neither is a memory hierarchy.

Retrieval-augmented generation (RAG) introduces external
memory but does not manage the context window itself. The retrieval
decision is made by an external system; the context window remains
an append-only buffer for retrieved content.

Recursive summarization (MemGPT~\cite{packer2023memgpt}) is the
closest prior work. MemGPT implements a two-tier memory with a main
context and an external archival store, using function calls to manage
paging. Our approach differs in three ways: (1) the projector is a
separate model from the reasoning model, enforcing separation of
concerns; (2) declared losses make every projection honest about what
was dropped; (3) the tensor is a structured representation with
navigable strands, not an unstructured summary.

Du et al.~\cite{du2025context} showed that context length alone
degrades LLM performance, with accuracy falling as context grows even
when the relevant information is present. This motivates the projector
architecture: the goal is not to retain more tokens but to retain
better-structured tokens.

The dispersion methodology is, to our knowledge, novel. Prior
evaluation of context representations relies on downstream task
accuracy or human judgment. Using an impossible task to isolate
representation quality from task performance has no direct precedent
in the literature.

\section{Discussion}

\subsection{The Append-Only Log Is the Wrong Architecture}

The deep raw log (Condition D) is the most expensive representation
and the worst performer. At 32,691 tokens, it consumes 7$\times$ the
budget of the tensor and produces 3$\times$ the dispersion. This is
not a marginal finding. The architecture that every production LLM
uses---append tokens, compact when full---is demonstrably inferior
to structured projection on the coherence metric.

The reason is architectural, not quantitative. An append-only log
preserves temporal ordering but not semantic structure. As the
log grows, signal-to-noise ratio decreases: early context that shaped
the reasoning trajectory is buried under later tool outputs,
corrections, and tangential exchanges. The model must reconstruct
semantic structure from temporal ordering at every invocation. The
tensor provides that structure directly.

\subsection{Structure Matters at Write Time, Not Read Time}

The ablation study's most surprising result is that
\texttt{content\_only}---unstructured prose---slightly outperforms the
full structured tensor at read time. The JSON scaffolding may impose a
parsing cost that partially offsets the navigational benefit.

But \texttt{content\_only} could not be produced without the tensor
structure at write time. The projector uses strands to organize its
integration, declared losses to track what was dropped, and branch
prediction to orient the projection. The structure guides the
projector's decisions even if the consumer does not need it. This
suggests an optimization: project with full structure, deliver as
prose. The memory controller uses a rich internal representation; the
ALU receives a simplified view.

\subsection{Meta-Cognition as Canary---and as Breath}

The precursor signal---declared losses or branch prediction going
empty---has a dual nature. In approximately 20\% of occurrences, it
precedes content collapse: the projector loses its ability to reason
about its own state before it loses the state itself. In hardware terms,
the MMU fails before the memory fails. In the remaining $\sim$80\%, the
signal is \emph{functional}: the projector temporarily sheds
meta-cognitive overhead to restructure its content representation,
then regenerates meta-cognition from the restructured content.

A trajectory ablation experiment (50 cycles $\times$ 4 conditions)
reveals the mechanism. With full meta-cognitive feedback (the tensor
receives its own prior declared losses), meta-cognitive allocation
\emph{increases} over time (slope $+$0.004) and the breathing rhythm
persists. Without loss feedback, meta-cognitive allocation atrophies
(slope $-$0.003), the breathing rhythm dampens, and loss generation
itself declines---the mirror creates the thing it reflects. Without
branch prediction feedback, the model compensates by generating more
loss declarations, but overall meta-cognitive capacity still declines.

This suggests that declared losses serve a different function than
branch prediction. Branch prediction (\texttt{instructions\_for\_next})
is \emph{outcome-load-bearing}: removing it degrades Riemann dispersion
by 22\%. Declared losses are \emph{process-load-bearing}: removing them
does not change outcome quality but changes the trajectory dynamics---
specifically, the tensor's ability to maintain meta-cognitive capacity
through the breathing/defragmentation cycle. The loss declarations
function as a recovery template: after a breathing event, the model
reconstructs its meta-cognitive scaffolding using the pattern
established by prior loss declarations.

The practical implication remains: monitor for precursor events. But
the correct response is to escalate on \emph{sustained} precursors
(two or more consecutive cycles), not individual events. Single-cycle
precursors are healthy breathing, not warning signs.

\section{Limitations}

The evaluation uses a single source session and a single reasoning
model (Claude Sonnet 4.6). Cross-session and cross-model replication
is needed to establish generality. The Riemann dispersion methodology
measures coherence of failure, not quality of success; it is possible
that a representation with lower dispersion on an impossible task does
not produce better outputs on feasible tasks. The growth curves run on
Haiku 4.5 only; the collapse ceiling likely differs for larger models.
The escalation experiment is not yet evaluated with the Riemann
methodology. The ablation study uses N=20, which may be insufficient
to detect small effects.

\section{Conclusion}

The append-only log with compaction is not the only architecture for
LLM context management. Structured tensor projection---with a
separate memory controller, declared losses, strand organization, and
collapse detection---produces tighter cognitive coherence than either
raw accumulation or production compaction. The dispersion methodology
provides a judge-free, rubric-free way to measure this. The growth
dynamics reveal a characteristic sawtooth pattern with a detectable
precursor signal. Together, these results suggest that LLM context
management would benefit from the same architectural thinking that
transformed flat-memory computers into hierarchical memory systems.

\bibliographystyle{plain}
\bibliography{references}

\end{document}
